\chapter{CONSIDERAÇÕES PARCIAIS E PRÓXIMAS ETAPAS}
\label{cap:consideracoes}

\section{Situação das atividades}
\label{sec:situacao}

O Quadro~\ref{qua:situacao} compara as atividades previstas na proposta com o estado atual. A parte de software está praticamente pronta. Faltam a calibração da proximidade, a montagem física do bracelete e as atividades com o público-alvo.

\begin{quadro}[htb]
\caption{Situação das atividades previstas na proposta}
\label{qua:situacao}
\small
\begin{tabular}{|p{5.4cm}|p{2.4cm}|p{6.2cm}|}
\hline
\textbf{Atividade prevista} & \textbf{Situação} & \textbf{Observação} \\
\hline
1. Pesquisa e levantamento de requisitos & Concluída & Ampliada para BLE, RSSI e Web Bluetooth após a mudança de escopo \\
\hline
2. Projeto e montagem do bracelete & Parcial & Projeto definido; montagem de motor, bateria e pulseira pendente \\
\hline
3. Programação do ESP32 e integração & Concluída & Firmware do bracelete e da baliza, interface web e modo demonstração \\
\hline
4. Testes, avaliação e apresentação & Em andamento & Testes funcionais feitos; calibração, testes com usuários e apresentação pendentes \\
\hline
5. Documentação e registro dos resultados & Em andamento & Código versionado e este relatório parcial \\
\hline
\end{tabular}
\normalsize
\fonte{Elaborado pelos autores (2026).}
\end{quadro}

\section{Considerações parciais}
\label{sec:consideracoes-parciais}

Até aqui, a ação mostrou que é possível construir, com duas placas de baixo custo e um celular comum, um sistema em que o bracelete detecta vários pontos de referência e os informa ao usuário por vibração e voz, sem aplicativo instalado e com comunicação em sentido único. A mudança de escopo, de detecção de obstáculos para orientação por balizas, manteve o objetivo geral da proposta, que é ampliar a percepção do ambiente por pessoas com baixa visão. Ela, porém, traz uma limitação que deve ficar clara: o sistema depende de o ambiente ter sido preparado com balizas e não detecta obstáculos não sinalizados.

O principal aprendizado técnico desta etapa foi a distância entre o modelo teórico de propagação e o comportamento real do hardware. A estimativa em metros, que parecia simples na Equação~\ref{eq:distancia}, mostrou-se pouco confiável com antenas pequenas, e a decisão de trocar metros por zonas veio diretamente dos testes. Outro aprendizado foi o cuidado com a informação dada ao usuário. Em uma tecnologia assistiva, anunciar ``caminho livre'' sem ter leitura, ou uma distância errada, pode ser pior do que não anunciar nada.

Para a formação dos extensionistas, a etapa exigiu integrar conteúdos de sistemas embarcados, redes sem fio, programação web e acessibilidade em um mesmo problema, além de lidar com questões práticas que não aparecem em sala de aula, como diferenças entre variantes de microcontrolador, pinos de inicialização e limitações de antena.

\section{Próximas etapas}
\label{sec:proximas}

\begin{enumerate}
    \item \textbf{Calibração:} medir o RSSI bruto com as placas encostadas, a 1~m e a distâncias maiores, em várias orientações, para definir os limites das zonas de proximidade;
    \item \textbf{Zonas de proximidade:} substituir os metros por zonas (``muito perto'', ``perto'' e ``na sala'') no firmware e na interface;
    \item \textbf{Montagem física:} ligar o motor de vibração com transistor de acionamento, alimentar por bateria LiPo e acomodar o conjunto em uma pulseira;
    \item \textbf{Novas balizas:} gravar ou adquirir balizas para, no mínimo, três pontos de uma sala do campus;
    \item \textbf{Validação com o público-alvo:} realizar testes supervisionados com pessoas com baixa visão, colhendo sua percepção sobre a clareza dos alertas;
    \item \textbf{Apresentação e relatório final:} apresentar o protótipo à comunidade acadêmica do IFPR -- Campus Pinhais e registrar os resultados no relatório final.
\end{enumerate}
